\begin{lemma}
Let $A$ be finite, $0\le p\le1$, $K\subseteq A$, and $q:A\to\mathbb R$
with $0\le q_a\le1$ for $a\in K$. On the Borel space $\mathbb R^2$ put
$\rho=(p\delta_1+(1-p)\delta_0)\otimes(\operatorname{Leb}|_{[0,1]})$,
and on $(\mathbb R^2)^{\{*\}\sqcup A}$ put $\mu=\rho^{\{*\}\sqcup A}$.
For $f_i=(b_i,t_i)$ define
\[
 L=\{f:\exists a\in K,\ \neg(b_*=1\land q_a<t_*\land t_*\in[0,1])\},
\quad
 R=\{f:\exists a\in K,\ \neg(b_a=1\land q_a<t_a\land t_a\in[0,1])\}.
\]
Both events are measurable and $\mu$ has total mass one. If $K$ is empty,
$\mu(L)=0$. For every maximizing $a\in K$, namely one with
$q_b\le q_a$ for all $b\in K$,
\[
 \mu(L)=\operatorname{ofReal}(1-p(1-q_a)).
\]
In all cases,
\[
 \mu(R)=\operatorname{ofReal}\left(1-\prod_{a\in K}p(1-q_a)\right),
 \qquad \mu(L)\le\mu(R).
\]
\end{lemma}
\begin{proof}
All coordinate evaluations are Borel measurable. Each failure test is the
intersection of a closed activation equality, an open strict priority
inequality, and a closed interval condition. Taking complements and finite
unions proves the stated measurability.

The two Dirac masses are at distinct points, so the activation measure has
mass $p$ at $1$, mass $1-p$ at $0$, and total mass one. Lebesgue measure of
$[0,1]$ is one. Thus $\rho$, and every finite product of copies of $\rho$,
has total mass one, including an empty product. For $0\le q\le1$, the
failure test on one coordinate is exactly $\{1\}\times(q,1]$ and has
mass $p(1-q)$ by the rectangle formula for product measure and the length
formula for intervals. Its complement therefore has mass $1-p(1-q)$.
These statements also hold at $p=0,1$ and $q=0,1$; no division has occurred.

If $K=\varnothing$, both displayed existential events are empty and the
empty product in the formula for $R$ is one. Otherwise a maximum of the
finite nonempty set of priorities is attained, say at $a$. Simultaneous
failure of all candidates under the common coordinate is equivalent to
$b_*=1$ and $q_a<t_*\le1$. Indeed this condition implies every failure,
and failure for the maximizing candidate implies this condition. The
common-coordinate success event consequently has mass $1-p(1-q_a)$.
Every unused coordinate contributes its total mass one.

The complement of $R$ requires the failure rectangle on coordinate $a$
for each $a\in K$, and imposes no restriction on any other coordinate.
The finite product rectangle formula gives its mass
$\prod_{a\in K}p(1-q_a)$. Subtracting from the unit total mass proves
the formula for $R$. This subtraction is finite; all masses lie in $[0,1]$.
Each factor lies in $[0,1]$. For a maximizing $a$, split the product into
its $a$ factor and the remaining product. The latter is at most one by
induction on its finite index set, so the whole product is at most
$p(1-q_a)$. Subtract from one to obtain $\mu(L)\le\mu(R)$.
The empty case was already proved. Applying $\operatorname{ofReal}$
preserves these equalities and inequalities between nonnegative reals.
\end{proof}
